\chapter{WebSocket, REST, TLS and JSON at Speed}\label{ch:ll:websocket-rest-tls-and-json-at-speed}

Crypto venues and many brokers do not speak FIX or binary protocols to most of their clients: market data arrives as JSON
text inside \omterm{def:ll:websocket-rest-tls-and-json-at-speed:ws}{WebSocket frames} inside TLS records, and orders leave as signed HTTP requests or as JSON over the same kind of
connection. A trading server placed in the venue's cloud region (One Quant Book~3, chapter~26) sits a fraction of a
millisecond from its matching engine. On this book's laptop, the Python reference of this chapter's build spends almost
\qty{50}{\micro\second} of processor time to receive one depth update and send one order (opening and sealing TLS records,
unmasking and masking frames, decoding the JSON, signing); the C++ components spend about one. The first is a tenth of a
half-millisecond round trip; the second disappears in it. This chapter takes the web stack apart layer by layer: HTTP and
the cost of opening connections, TLS and what its handshake and records cost, \omterm{def:ll:websocket-rest-tls-and-json-at-speed:ws}{WebSocket} framing, JSON decoding with a
structural index, and request signatures computed without rehashing the key.

\section{HTTP, REST and connection reuse}

\begin{definition}[REST interface, persistent connection, connection pool]\label{def:ll:websocket-rest-tls-and-json-at-speed:rest}
A \emph{REST interface}\index{REST interface} exposes a service as resources addressed by URLs and manipulated by HTTP
requests (GET to read, POST to create, DELETE to remove), each request carrying everything the server needs to process it. A
\emph{persistent connection}\index{persistent connection} carries many HTTP requests and responses over one TCP (and TLS)
connection instead of opening one per request. A \emph{connection pool}\index{connection pool} keeps a set of open persistent
connections to a server and lends an idle one to each request.
\end{definition}

A crypto venue's order entry is typically a \omterm{def:ll:websocket-rest-tls-and-json-at-speed:rest}{REST interface}: an order is a POST with the order's parameters, the client's API
key and a signature (One Quant Book~3, chapter~15), and the answer is JSON. The expensive part is not the request but the
connection under it. Opening one costs a TCP handshake, one round trip, then a \omterm{def:ll:websocket-rest-tls-and-json-at-speed:tls}{TLS handshake}, another round trip plus the
cryptography of key exchange and certificate verification; a \omterm{def:ll:websocket-rest-tls-and-json-at-speed:ws}{WebSocket} adds an HTTP upgrade, a third. HTTP/1.1 ``defaults to
the use of \omterm{def:ll:websocket-rest-tls-and-json-at-speed:rest}{persistent connections}'' (RFC~9112), and a client that uses them pays these costs once per connection instead of
once per order. The build's pool keeps a fixed number of connections open and passes every request through Book~3's rate-limit
governor first, since a venue's limits count requests and weights (Book~3, chapter~15), not connections.

\begin{omfigure}
\begin{tikzpicture}[font=\scriptsize,>=Stealth]
\def\L{0}\def\R{7.5}
\node[font=\footnotesize\bfseries] at (\L,0.45) {client};
\node[font=\footnotesize\bfseries] at (\R,0.45) {venue};
\draw[gray,thick] (\L,0.15) -- (\L,-8.1);
\draw[gray,thick] (\R,0.15) -- (\R,-8.1);
\newcommand{\msgR}[3]{\draw[->,#3] (\L,#1) -- node[above,sloped]{#2} (\R,#1-0.3);}
\newcommand{\msgL}[3]{\draw[<-,#3] (\L,#1-0.3) -- node[above,sloped]{#2} (\R,#1);}
\msgR{-0.1}{TCP SYN}{gray}
\msgL{-0.8}{SYN-ACK}{gray}
\msgR{-1.5}{ACK, TLS ClientHello with a key share}{omDef}
\msgL{-2.2}{ServerHello, certificate, Finished}{omDef}
\msgR{-2.9}{Finished, HTTP GET with Upgrade: websocket}{omMeth}
\msgL{-3.6}{101 Switching Protocols}{omMeth}
\msgL{-4.5}{depth update: TLS record [WebSocket frame [JSON]]}{omThm,thick}
\msgR{-5.2}{order: TLS record [masked frame [JSON, signature]]}{omThm,thick}
\msgL{-6.4}{ping (every 20 s at one venue)}{gray}
\msgR{-7.1}{pong}{gray}
\draw[decorate,decoration={brace,amplitude=4pt,mirror},thick] (-0.25,-0.1) -- (-0.25,-3.9)
  node[midway,left=5pt,align=right]{three round trips\\ before the first\\ message};
\draw[decorate,decoration={brace,amplitude=4pt},thick] (7.75,-4.5) -- (7.75,-5.5)
  node[midway,right=5pt,align=left]{on an open\\ connection: one\\ record each way};
\end{tikzpicture}

\omcaption{A \omterm{def:ll:websocket-rest-tls-and-json-at-speed:ws}{WebSocket} connection over TLS~1.3: one round trip for TCP, one for TLS, one for the HTTP upgrade; then each
message is a \omterm{def:ll:websocket-rest-tls-and-json-at-speed:ws}{WebSocket frame} inside a TLS record. A REST order on a new connection pays the first two round trips; on a
\omterm{def:ll:websocket-rest-tls-and-json-at-speed:rest}{persistent connection}, none.}\label{fig:ll:websocket-rest-tls-and-json-at-speed:setup}
\end{omfigure}

\section{TLS: the handshake and the record layer}

\begin{definition}[TLS handshake, session resumption]\label{def:ll:websocket-rest-tls-and-json-at-speed:tls}
The \emph{TLS handshake}\index{TLS handshake} is the exchange that opens a TLS connection: the two sides agree on a protocol
version and cipher, exchange key shares, the server proves its identity with its certificate, and both derive the keys that
protect the records that follow; in TLS~1.3 it takes one round trip. \emph{Session resumption}\index{session resumption}
opens a new connection with a key derived from an earlier one (a pre-shared key sent to the client as a ticket), skipping
the certificate and part of the key exchange.
\end{definition}

TLS~1.3 (RFC~8446) cut the full handshake to one round trip and replaced the older session identifiers and tickets by
resumption with a pre-shared key: ``the key derived from the initial handshake is used to bootstrap the cryptographic state
instead of a full handshake''. A zero-round-trip mode lets a resuming client send data with its first message, ``at the cost
of certain security properties'' (the data can be replayed), which is why an order must never travel in it. After the
handshake, the record layer cuts the stream into records of at most $2^{14}$ bytes and protects each with an authenticated
cipher.

\Cref{tab:ll:websocket-rest-tls-and-json-at-speed:tls} gives the costs on the laptop's loopback, where the network costs
almost nothing and what remains is processing. A full TLS~1.3 handshake with an elliptic-curve certificate takes about
\qty{2.1}{\milli\second} with Python's \texttt{ssl} module (both ends on the same machine), a resumed one about
\qty{1.6}{\milli\second}. A round trip of 300 bytes on an open connection takes about \qty{45}{\micro\second} over TLS and
\qty{33}{\micro\second} over plain TCP, both in Python; the protection itself, one AES-128-GCM seal or open of a 300-byte
record with OpenSSL called from C++, is under \qty{0.2}{\micro\second}. The handshake is ten thousand times dearer than a record:
the reason to keep connections open.

\begin{table}[ht]
\centering\small
\begin{tabular}{lr}
\toprule
Operation (loopback, laptop) & Cost \\
\midrule
TLS~1.3 full handshake (Python \texttt{ssl}, both ends) & \qty{2.1}{\milli\second} \\
TLS~1.3 resumed handshake (Python \texttt{ssl}, both ends) & \qty{1.6}{\milli\second} \\
Round trip of 300 bytes over TLS (Python) & \qty{45}{\micro\second} \\
Round trip of 300 bytes over plain TCP (Python) & \qty{33}{\micro\second} \\
AES-128-GCM seal of a 300-byte record (OpenSSL, C++) & \qty{0.16}{\micro\second} \\
AES-128-GCM open of a 300-byte record (OpenSSL, C++) & \qty{0.15}{\micro\second} \\
\bottomrule
\end{tabular}
\caption{TLS on the laptop's loopback: medians of 300 handshakes each way and of 5\,000 round trips, and of 21 passes of
1\,000 records. Measured on a laptop (Intel Core Ultra 7 155H) under WSL2, no isolated cores; OpenSSL~3.0.2, Python~3.10.
Data: \texttt{bench\_web.py}, \texttt{measured\_tls.csv} and \texttt{measured\_web.csv}.}\label{tab:ll:websocket-rest-tls-and-json-at-speed:tls}
\end{table}

\section{WebSocket framing}

\begin{definition}[WebSocket, WebSocket frame]\label{def:ll:websocket-rest-tls-and-json-at-speed:ws}
A \emph{WebSocket}\index{WebSocket} is a bidirectional message channel over one TCP connection (usually inside TLS), opened by
an HTTP request that asks to upgrade the connection (RFC~6455). A \emph{WebSocket frame}\index{WebSocket frame} is its unit
on the wire: a two-byte header (a final-fragment bit, an opcode for text, binary, continuation, close, ping or pong, a mask
bit and a 7-bit length), an extended length of 2 or 8 bytes when the payload exceeds 125 bytes, a 4-byte masking key on frames
from the client, and the payload.
\end{definition}

Three rules of RFC~6455 shape a client. ``A client MUST mask all frames that it sends to the server'', with a new random key
per frame, even over TLS: the payload is XORed with the key repeated, a defence of intercepting proxies, not of the data. A
message may be split into fragments (a first frame with the opcode, continuation frames, the last with its final bit), and
control frames (ping, pong, close) may arrive between fragments, carry at most 125 bytes and are never fragmented. And the
venue checks the connection with pings that the client must answer: one large venue sends a ping every 20 seconds and drops
a connection that has not answered within a minute. The build's decoder unmasks in place eight bytes at a time.

\omcode{code/firm/wsclient/cpp/firm_wsclient.hpp}{26}{39}{Unmasking a payload with the 4-byte key repeated in a 64-bit
word.}\label{lst:ll:websocket-rest-tls-and-json-at-speed:unmask}

Framing costs little in C++: about \qty{24}{\nano\second} to decode and unmask a 300-byte frame, \qty{16}{\nano\second} to
encode and mask one. The build's Python reference, which XORs the payload one byte at a time, takes about
\qty{16}{\micro\second} for the same frame: several hundred times more, and the largest item of its budget.

\section{JSON at speed}

\begin{definition}[Structural index]\label{def:ll:websocket-rest-tls-and-json-at-speed:index}
A \emph{structural index} of a JSON text is the array of the positions of its structural characters (braces, brackets,
colons, commas) and quotes, computed in one vectorised pass; a decoder then walks the index instead of the bytes, reading each
value directly between two known positions.
\end{definition}

The parser of Langdale and Lemire (chapter~14) works in two stages: the first writes ``the location of all structural
characters \dots\ as integer indexes in a separate array'', the second walks that index. The build applies the idea to one
message shape, the depth update a large venue documents for its diff-depth stream: an event time, a symbol, the first and last
update identifiers (\texttt{U}, \texttt{u}) that Book~3's book builder checks for gaps, and arrays of bids and asks as
\texttt{[price, quantity]} strings. The first stage is chapter~14's \texttt{positions\_any} over the seven characters
\texttt{\{\}[]:,} and the double quote; the second walks the positions with the shape known in advance, reads each number where it stands, and
parses prices into integers of $10^{-8}$ without floating point. Anything unexpected (an escape, a missing key) makes it
return false, and a general parser takes over.

\omcode{code/firm/wsclient/cpp/firm_wsclient.hpp}{177}{204}{The second stage: walking the structural index of a depth
update.}\label{lst:ll:websocket-rest-tls-and-json-at-speed:walk}

On the fixture's 1\,000 depth updates (about 286 bytes each, up to ten levels a side), the indexed decoder takes about
\qty{0.39}{\micro\second} per update; a general parser that builds a tree of maps, vectors and strings and converts prices
with \texttt{strtod} about \qty{2.6}{\micro\second}, seven times more; Python's \texttt{json} module and the conversion to
integers about \qty{10}{\micro\second}. The decoded levels, divided by the instrument's tick and lot, are what Book~3's
\texttt{firm.wsbook} applies.

\section{Signing requests off the hot path}

\begin{definition}[Request signature]\label{def:ll:websocket-rest-tls-and-json-at-speed:signature}
A \emph{request signature} is a message authentication code computed by the client over a request's parameters with a secret
shared with the venue (typically HMAC-SHA-256 of the query string and body, sent in hexadecimal with the API key), which lets
the venue check that the request comes from the key's owner and was not altered; a timestamp and a validity window inside the
signed payload stop old requests from being replayed.
\end{definition}

HMAC (RFC~2104) computes $H\big((K \oplus \mathit{opad}) \,\|\, H((K \oplus \mathit{ipad}) \,\|\, m)\big)$. Both padded keys
are one 64-byte block, so the hash's state after absorbing each of them depends only on the key; the RFC notes that these
intermediate results ``can be precomputed only once at the time of generation of the key''. A signer that does so absorbs the
key once, copies the two states for each request and hashes only the message. The build's signer does this in C++, Rust and
Python, and is checked against RFC~4231's vectors and against the worked example of a large venue's documentation, whose
published key and order payload give a published signature.

\omcode{code/firm/wsclient/cpp/firm_wsclient.hpp}{279}{301}{The key's pads hashed once, at construction; each signature
copies two states.}\label{lst:ll:websocket-rest-tls-and-json-at-speed:hmac}

The costs teach two lessons. Precomputing the pads saves about two-fifths with the build's portable SHA-256 (about
\qty{0.64}{\micro\second} against \qty{1.1}{\micro\second} for a 119-byte order), and four-fifths with OpenSSL (about
\qty{0.19}{\micro\second} against \qty{0.92}{\micro\second} for its one-shot \texttt{HMAC()}). And the library is several times
faster than a portable implementation because it uses the processor's SHA instructions: cryptography is the one place where
the fast path is someone else's code, called well.

\begin{dated}{2026-09}{A venue's signing and streams}\label{dat:ll:websocket-rest-tls-and-json-at-speed:venue}
Binance's spot API documentation (GitHub, consulted September 2026) signs each request with HMAC-SHA-256 of the query string
and body, keyed by the API key's secret, and bounds its validity with a timestamp and a window, \texttt{recvWindow}, of 5\,000
milliseconds by default and 60\,000 at most. Its \omterm{def:ll:websocket-rest-tls-and-json-at-speed:ws}{WebSocket} streams send a ping every 20 seconds, drop a connection that has not
answered within a minute, and end any connection after 24 hours; its diff-depth stream publishes updates every 100 or 1\,000
milliseconds.
\end{dated}

\Cref{fig:ll:websocket-rest-tls-and-json-at-speed:budget} puts the pieces together: the processor time of receiving one depth
update and sending one order, step by step, for the build's C++ components and for its Python reference.

\begin{omfigure}
\begin{tikzpicture}
\begin{axis}[width=12cm,height=5.6cm,ybar,bar width=9pt,ymode=log,log origin y=infty,ymin=0.01,ymax=100,
  symbolic x coords={tls_open,frame_decode,json,sign,frame_encode,tls_seal},xtick=data,
  xticklabels={{open TLS\\ record},{decode\\ frame},{decode\\ JSON},{sign\\ order},{encode and\\ mask frame},{seal TLS\\ record}},
  xticklabel style={align=center,font=\footnotesize},enlarge x limits=0.1,
  ylabel={\si{\micro\second} (log scale)},tick label style={font=\small},label style={font=\small},grid=major,
  grid style={gray!25},legend columns=2,legend style={font=\footnotesize,at={(0.5,-0.3)},anchor=north,draw=none,fill=none,
  /tikz/every even column/.append style={column sep=8pt}},table/col sep=comma]
\addplot[fill=omDef!70,draw=omDef] table[x=step,y=cpp]{figdata/low-latency/17-websocket-rest-tls-and-json-at-speed/measured_budget_steps.csv};
\addplot[fill=omThm!45,draw=omThm] table[x=step,y=python]{figdata/low-latency/17-websocket-rest-tls-and-json-at-speed/measured_budget_steps.csv};
\legend{C++ components (0.9~\si{\micro\second} in all),Python reference (48~\si{\micro\second} in all)}
\end{axis}
\end{tikzpicture}

\omcaption{The processor time of one websocket order: receiving a depth update (TLS record, frame, JSON) and sending an order
(signature, masked frame, TLS record). C++: the build's decoders and OpenSSL; Python: the build's reference and \texttt{ssl}
(its record cost estimated as a quarter of the extra round-trip time over TLS). Measured on a laptop (Intel Core Ultra 7 155H)
under WSL2, no isolated cores. Data: \texttt{bench\_web.py}.}\label{fig:ll:websocket-rest-tls-and-json-at-speed:budget}
\end{omfigure}

\section{Tutorial: the web stack, measured}\label{tut:ll:websocket-rest-tls-and-json-at-speed:tut}

\begin{tutorial}
\textbf{Goal.} Build and check each layer, then measure the cost of one websocket order. \textbf{End state:}
\cref{tab:ll:websocket-rest-tls-and-json-at-speed:tls}, \cref{fig:ll:websocket-rest-tls-and-json-at-speed:budget}, and
green tests in Python, C++ and Rust.
\begin{tutsteps}
\item \textbf{Frames.} The tests encode RFC~6455's own examples (an unmasked and a masked ``Hello'', lengths needing 16 and 64
bits) and decode the fixture stream written by \texttt{make\_ws\_fixtures.py}: fifty depth updates, some fragmented in three
with a ping between the fragments, and a close. Python, C++ and Rust deliver the same messages in the same order.
\item \textbf{Depth updates.} The C++ indexed decoder reproduces, level for level, the Python reference's decoding of the
fixture's 1\,000 updates.
\item \textbf{Signatures.} All three signers reproduce RFC~4231's test cases and the venue's published example.
\item \textbf{TLS.} \texttt{bench\_web.py} makes a self-signed certificate with the \texttt{openssl} command, runs a TLS~1.3
server in a thread, and times full and resumed handshakes (it checks that every resumption was accepted) and round trips.
\end{tutsteps}
\textbf{What to change next.} Unmask with 32-byte vector XORs (chapter~14) and measure a 16-kilobyte frame; replace the Python
reference's masking with a bytes-level XOR over integers and measure again.
\end{tutorial}

\section{Build: the web client}\label{bld:ll:websocket-rest-tls-and-json-at-speed:wsclient}

\begin{build}
\bfield{Purpose} The firm's connection to venues that speak JSON over \omterm{def:ll:websocket-rest-tls-and-json-at-speed:ws}{WebSocket} and REST: depth updates decoded for Book~3's
\texttt{firm.wsbook}, orders signed and sent through a pool of \omterm{def:ll:websocket-rest-tls-and-json-at-speed:rest}{persistent connections} that respects the venue's rate limits.
\bfield{Interface} Python reference \texttt{firm\_wsclient}: \texttt{encode\_frame}, \texttt{decode\_frame},
\texttt{Reassembler}, \texttt{decode\_depth}, \texttt{to\_book\_levels}, \texttt{Signer}, \texttt{Pool}. C++20
\texttt{firm::ws}: \texttt{encode\_frame}, \texttt{decode\_frame} (unmasking in place), \texttt{Reassembler},
\texttt{decode\_depth} into a fixed \texttt{Depth}, \texttt{Sha256}, \texttt{HmacSha256}. Rust \texttt{firm\_wsclient}:
\texttt{encode\_frame}, \texttt{decode\_frame}, \texttt{Sha256}, \texttt{HmacSha256}.
\bfield{Rules} Client frames are always masked; control frames are never fragmented and carry at most 125 bytes; lengths are
minimally encoded; prices and quantities are integers of $10^{-8}$, never floats; the key's pads are hashed once; the pool
never opens more than its size and asks the governor before every request.
\bfield{Acceptance tests} \texttt{code/firm/wsclient/}: RFC~6455's frame examples; the fixture stream reassembled identically
in the three languages; the 1\,000 fixture updates decoded identically in Python and C++; RFC~4231 test cases 1, 2 and 6 and
the venue's example signature in the three languages; the pool's reuse, capacity and throttling.
\bfield{Stretch} Vectorised unmasking; a general structural-index JSON reader; per-message deflate; the order path through
the simulator's crypto preset of One Quant Book~10.
\end{build}

\begin{omsources}
\item IETF RFC~6455 (WebSocket), RFC~8446 (TLS~1.3), RFC~2104 (HMAC), RFC~4231 (HMAC-SHA-256 test vectors), RFC~9112
(HTTP/1.1).
\item G.~Langdale and D.~Lemire, ``Parsing gigabytes of JSON per second'', \emph{The VLDB Journal} 28(6), 2019.
\item Binance spot API documentation (GitHub): REST request security, WebSocket streams.
\end{omsources}

\section{Exercises}

\begin{exercise}[$\star$]\label{exo:ll:websocket-rest-tls-and-json-at-speed:1}
How many bytes does a client frame carrying 300 bytes of payload occupy? And 70\,000 bytes? And a server frame of 100 bytes?
\end{exercise}

\begin{exercise}[$\star$]\label{exo:ll:websocket-rest-tls-and-json-at-speed:2}
A REST order on a new connection pays TCP and TLS setup. With a round trip of \qty{0.5}{\milli\second} to the venue and the
measured handshake, how much later does it arrive than on a \omterm{def:ll:websocket-rest-tls-and-json-at-speed:rest}{persistent connection}?
\end{exercise}

\begin{exercise}[$\star$]\label{exo:ll:websocket-rest-tls-and-json-at-speed:3}
Why must a client never send an order in TLS~1.3's zero-round-trip data?
\end{exercise}

\begin{exercise}[$\star\star$]\label{exo:ll:websocket-rest-tls-and-json-at-speed:4}
Why can the two padded keys of HMAC be hashed in advance, and how many SHA-256 compressions does signing a 119-byte payload
take with and without that precomputation?
\end{exercise}

\begin{exercise}[$\star\star$]\label{exo:ll:websocket-rest-tls-and-json-at-speed:5}
A depth update's price is the string \texttt{64123.2}. What does the indexed decoder return at $10^{-8}$, and what can go wrong if a
decoder converts it with \texttt{strtod} and multiplies by $10^8$?
\end{exercise}

\begin{exercise}[$\star\star$]\label{exo:ll:websocket-rest-tls-and-json-at-speed:6}
A ping arrives between the first and second fragments of a depth update. What must the client do, and in what order?
\end{exercise}

\begin{exercise}[$\star\star\star$]\label{exo:ll:websocket-rest-tls-and-json-at-speed:7}
\emph{Coding.} Replace the Python reference's byte-by-byte masking with one XOR of two large integers built with
\texttt{int.from\_bytes}, check the fixture tests, and measure the frame costs again.
\end{exercise}

\begin{exercise}[$\star\star\star$]\label{exo:ll:websocket-rest-tls-and-json-at-speed:8}
\emph{Find the flaw.} ``Each strategy thread opens its own HTTPS connection for each order, so that no two threads ever
contend for a connection.''
\end{exercise}

\section{Problem: Cloud to Venue}

\begin{problem}[{Weekend problem --- where the time of a websocket order goes}]\label{pb:ll:websocket-rest-tls-and-json-at-speed:1}
A strategy runs in the same cloud region as a crypto venue. Assume a one-way network hop of \qty{250}{\micro\second} and a
decision of negligible cost; take the measured costs of \cref{fig:ll:websocket-rest-tls-and-json-at-speed:budget} and
\cref{tab:ll:websocket-rest-tls-and-json-at-speed:tls} (\texttt{measured\_budget.csv}, \texttt{measured\_tls.csv}).

\medskip\noindent\textbf{Part I --- The budget.}
\begin{enumerate}
\item What is the total processor time of one order with the C++ components, and with the Python reference?
\item Which step dominates each?
\item What share of the path venue, host, venue is the host's processing, in each case?
\item What would the Python share be with the C++ framing alone (the rest unchanged)?
\end{enumerate}

\medskip\noindent\textbf{Part II --- Connections.}
\begin{enumerate}[resume]
\item What does opening a new connection for each REST order add, with a full handshake, and with a resumed one?
\item What does a new \omterm{def:ll:websocket-rest-tls-and-json-at-speed:ws}{WebSocket} add before its first message?
\item How many orders a second could one connection carry if each took a new connection?
\item What does the venue's 24-hour connection limit imply for a strategy that runs for days?
\end{enumerate}

\medskip\noindent\textbf{Part III --- Signatures and data.}
\begin{enumerate}[resume]
\item Why can the signature not be computed before the decision, and what can be?
\item At 100 depth updates a second per instrument over 50 instruments, how much of a core does decoding take with each
decoder?
\item What does the update identifier pair \texttt{(U, u)} let the client detect, and what does it do then?
\item Why is the timestamp inside the signed payload?
\end{enumerate}

\medskip\noindent\textbf{Part IV --- The verdict.}
\begin{enumerate}[resume]
\item State the \emph{named result}: the share of the tick-to-trade spent in decryption, framing, JSON and signing against the
network hop, and what connection reuse saves per order.
\item When does the Python reference suffice?
\item Where would you put the TLS work if the budget were tighter still?
\item What would you measure on the real server before trusting these laptop numbers?
\item Why does the network hop not appear in the processor budget, and why does it still matter?
\item What part of this stack does a venue's binary protocol remove?
\item How would you test the pool's behaviour when the venue answers 429?
\item In one sentence: what is fast about a web stack that is fast?
\end{enumerate}
\end{problem}

\section{Interview questions}

\begin{interviewq}[$\star$ \iqroles{developer}]\label{iq:ll:websocket-rest-tls-and-json-at-speed:1}
Why reuse HTTP connections for order entry? What does a new connection cost?
\end{interviewq}

\begin{interviewq}[$\star\star$ \iqroles{developer}]\label{iq:ll:websocket-rest-tls-and-json-at-speed:2}
Describe a \omterm{def:ll:websocket-rest-tls-and-json-at-speed:ws}{WebSocket frame}. Why are client frames masked, and what does masking cost?
\end{interviewq}

\begin{interviewq}[$\star\star$ \iqroles{developer}]\label{iq:ll:websocket-rest-tls-and-json-at-speed:3}
How would you parse a venue's JSON depth updates as fast as possible?
\end{interviewq}

\begin{interviewq}[$\star\star$ \iqroles{developer}]\label{iq:ll:websocket-rest-tls-and-json-at-speed:4}
How is an HMAC request signature computed, and how would you make signing cheap?
\end{interviewq}

\begin{interviewq}[$\star\star$ \iqroles{developer}]\label{iq:ll:websocket-rest-tls-and-json-at-speed:5}
What does TLS~1.3 \omterm{def:ll:websocket-rest-tls-and-json-at-speed:tls}{session resumption} save, and what does zero-round-trip data risk?
\end{interviewq}

\begin{interviewq}[$\star\star\star$ \iqroles{developer, researcher}]\label{iq:ll:websocket-rest-tls-and-json-at-speed:6}
Your crypto strategy's orders arrive 3~ms after the market data that triggered them, although the venue is 0.3~ms away. Where
do you look?
\end{interviewq}
